% Blueprint for the Monogenic Extensions formalization project
% Based on arXiv:2503.07846 -- Lemmas 3.1 and 3.2
%
% Authors: Bianca Viray, Bryan Boehnke, Grant Yang, George Peykanu, Tianshuo Wang
% UW Math AI Winter 2026

\chapter{Introduction}

This blueprint describes the formalization of Lemmas 3.1 and 3.2 from
\cite{monogenic}, which concern \emph{monogenic extensions} of local rings.
The main results are:

\begin{itemize}
\item \textbf{Lemma 3.2} (the étale case): A finite injective étale map of local rings
  is monogenic, and the minimal polynomial has unit derivative.
\item \textbf{Lemma 3.1} (the partially étale case): If a finite map of local rings
  admits an étale quotient through a height-one prime, the extension is monogenic.
\item \textbf{Converse}: A monogenic extension with unit derivative is étale
  (i.e., standard étale).
\end{itemize}

The formalization builds on Mathlib's theory of étale ring homomorphisms,
local rings, residue fields, and the primitive element theorem.

\begin{thebibliography}{9}
\bibitem{monogenic}
  B.~Viray et al.,
  \emph{Monogenic extensions of local rings},
  arXiv:2503.07846v2, 2025.
\end{thebibliography}


\chapter{The Quotient Isomorphism}
\label{ch:generator}

This chapter establishes that when $R[\beta] = S$ for a finite free extension,
the natural map $R[X]/(\operatorname{minpoly}_R \beta) \to S$ is an isomorphism,
\emph{without} assuming $R$ is integrally closed.

\section{Degree Bound via Cayley--Hamilton}

\begin{lemma}[Minimal polynomial degree bound]
  \label{lem:minpoly_natDegree_le}
  \lean{minpoly.natDegree_le'}
  \leanok
  Let $R \to S$ be a finite free extension with $R$ nontrivial.
  For any $\alpha \in S$,
  \[
    \deg(\operatorname{minpoly}_R \alpha) \le \operatorname{finrank}_R S.
  \]
\end{lemma}

\begin{proof}
  \leanok
  Let $\ell_\alpha = \operatorname{Algebra.lmul}_R\, \alpha$ be the
  left-multiplication endomorphism.  By Cayley--Hamilton
  ($\mathtt{LinearMap.aeval\_self\_charpoly}$),
  $\operatorname{aeval}_\alpha(\operatorname{charpoly}(\ell_\alpha)) = 0$.
  Since $\operatorname{charpoly}(\ell_\alpha)$ is monic, the minimal polynomial
  divides it, so
  $\deg(\operatorname{minpoly}_R \alpha) \le \deg(\operatorname{charpoly}(\ell_\alpha)) = \operatorname{finrank}_R S$.
  \end{proof}

\section{The Isomorphism via Orzech's Property}

\begin{theorem}[Quotient isomorphism without integrally closed hypothesis]
  \label{thm:mkOfAdjoinEqTop}
  \lean{IsAdjoinRootMonic.mkOfAdjoinEqTop'}
  \leanok
  \uses{lem:minpoly_natDegree_le}
  Let $R \to S$ be a finite free extension with $R$ nontrivial,
  and let $\alpha \in S$ with $R[\alpha] = S$.
  Then $f = \operatorname{minpoly}_R \alpha$ is monic, and the natural map
  \[
    \varphi \colon R[X]/(f) \longrightarrow S, \qquad X \mapsto \alpha
  \]
  is an $R$-algebra isomorphism. In particular, $S$ admits an
  $\mathtt{IsAdjoinRootMonic}$ structure for $f$.
\end{theorem}

\begin{proof}
  \leanok
  \uses{lem:minpoly_natDegree_le}
  The map $\varphi$ is well-defined since $f(\alpha) = 0$, and surjective
  since $R[\alpha] = S$.

  \textbf{Rank equality.}
  The upper bound $\deg f \le \operatorname{finrank}_R S$
  follows from Lemma~\ref{lem:minpoly_natDegree_le}.
  For the lower bound, consider the surjection $\varphi$.
  By the quotient rank formula
  ($\mathtt{finrank\_quotient\_}\allowbreak\mathtt{span\_eq\_natDegree'}$),
  $\operatorname{finrank}_R(R[X]/(f)) = \deg f$.
  Since $\varphi$ is surjective,
  $\operatorname{finrank}_R S \le \operatorname{finrank}_R(R[X]/(f)) = \deg f$.
  Hence $\deg f = \operatorname{finrank}_R S$.

  \textbf{Injectivity.}
  Since both $R[X]/(f)$ and $S$ are free $R$-modules of the same finite rank,
  and $\varphi$ is surjective, the composite
  $e^{-1} \circ \varphi$ is a surjective endomorphism of a finitely generated
  free module (where $e$ is a linear equivalence from the rank equality).
  By $\mathtt{OrzechProperty.}\allowbreak\mathtt{injective\_of\_surjective\_endomorphism}$
  (which holds for all commutative rings), this composite is injective,
  hence $\varphi$ is injective.
\end{proof}


\chapter{The \'Etale Case (Lemma 3.2)}
\label{ch:etale}

We now prove the main result for finite étale extensions of local rings:
such extensions are monogenic, and the minimal polynomial has unit derivative.

\section{Generator Descends to Residue Field}

\begin{lemma}[Generator descends to residue fields]
  \label{lem:adjoin_residue_eq_top}
  \lean{Monogenic.adjoin_residue_eq_top_of_adjoin_eq_top}
  \leanok
  If $\beta$ generates $S$ over $R$ (i.e., $R[\beta] = S$),
  then $\bar\beta = \beta \bmod \mathfrak{m}_S$ generates $k_S$ over $k_R$.
\end{lemma}

\begin{proof}
  \leanok
  Given $x \in k_S$, lift to $s \in S$, express $s = p(\beta)$ for some $p \in R[X]$,
  then $x = \bar{p}(\bar\beta)$ where $\bar{p} = p \bmod \mathfrak{m}_R \in k_R[X]$.
  The key commutativity ($\operatorname{residue}_S \circ \varphi
  = \operatorname{algebraMap}_{k_R \to k_S} \circ \operatorname{residue}_R$)
  is used via $\mathtt{map\_aeval\_eq\_aeval\_map}$.
\end{proof}

\section{Rank Equality across Residue Fields}

\begin{lemma}[Rank preserved by residue field base change]
  \label{lem:finrank_eq_finrank_residueField}
  \lean{Monogenic.finrank_eq_finrank_residueField}
  \leanok
  If $R \to S$ is a finite étale extension of local rings with faithful
  scalar action, then
  \[
    \operatorname{finrank}_R S = \operatorname{finrank}_{k_R} k_S.
  \]
\end{lemma}

\begin{proof}
  \leanok
  Since $R \to S$ is étale, it is smooth, hence flat.
  Finite + flat over a local ring implies free
  ($\mathtt{Module.free\_of\_flat\_of\_isLocalRing}$).

  The étale condition gives $\mathfrak{m}_R \cdot S = \mathfrak{m}_S$
  ($\mathtt{Algebra.FormallyUnramified.}\allowbreak\mathtt{map\_maximalIdeal}$), so
  $S / \mathfrak{m}_R S \cong k_S$.
  By $\mathtt{IsLocalRing.finrank\_quotient\_map}$,
  $\operatorname{finrank}_{k_R}(S/\mathfrak{m}_R S) = \operatorname{finrank}_R S$.
  The quotient equivalence from $\mathfrak{m}_R S = \mathfrak{m}_S$
  transfers the finrank.
\end{proof}

\section{Minimal Polynomial Descends to Residue Field}

\begin{lemma}[Minimal polynomial maps to residue minimal polynomial]
  \label{lem:minpoly_map_residue}
  \lean{Monogenic.minpoly_map_residue}
  \leanok
  \uses{lem:adjoin_residue_eq_top, lem:finrank_eq_finrank_residueField, thm:mkOfAdjoinEqTop}
  Let $R \to S$ be a finite étale extension of local rings with faithful
  scalar action, and $\beta \in S$ with $R[\beta] = S$.
  Then
  \[
    (\operatorname{minpoly}_R \beta) \bmod \mathfrak{m}_R
    = \operatorname{minpoly}_{k_R} \bar\beta.
  \]
\end{lemma}

\begin{proof}
  \leanok
  \uses{lem:adjoin_residue_eq_top, lem:finrank_eq_finrank_residueField, thm:mkOfAdjoinEqTop}
  Set $f = \operatorname{minpoly}_R \beta$, $\bar f = f \bmod \mathfrak{m}_R$,
  $g = \operatorname{minpoly}_{k_R} \bar\beta$.

  Since $\bar\beta$ is a root of $\bar f$ (by $\mathtt{hom\_eval_2}$),
  we have $g \mid \bar f$.
  Both $f$ and $g$ are monic, and $\bar f$ is monic.

  The degree chain:
  \begin{align*}
    \deg \bar f &= \deg f
      && \text{(monic polynomial preserves degree under reduction)} \\
    &= \operatorname{finrank}_R S
      && \text{(Theorem~\ref{thm:mkOfAdjoinEqTop})} \\
    &= \operatorname{finrank}_{k_R} k_S
      && \text{(Lemma~\ref{lem:finrank_eq_finrank_residueField})} \\
    &= [k_R(\bar\beta) : k_R]
      && \text{(Lemma~\ref{lem:adjoin_residue_eq_top}: } k_R[\bar\beta] = k_S \text{)} \\
    &= \deg g
      && \text{(degree of minimal polynomial = extension degree)}
  \end{align*}
  Since $g \mid \bar f$, both are monic, and they have the same degree,
  $\bar f = g$.
\end{proof}

\section{Unit Derivative Condition}

\begin{theorem}[Derivative of minimal polynomial is a unit]
  \label{thm:isUnit_aeval_derivative_minpoly}
  \lean{Monogenic.isUnit_aeval_derivative_minpoly_of_adjoin_eq_top}
  \leanok
  \uses{lem:minpoly_map_residue}
  Let $R \to S$ be a finite étale extension of local rings with faithful
  scalar action, and $\beta \in S$ with $R[\beta] = S$.  Then
  \[
    f'(\beta) \in S^\times
  \]
  where $f = \operatorname{minpoly}_R \beta$.
\end{theorem}

\begin{proof}
  \leanok
  \uses{lem:minpoly_map_residue}
  Since $S$ is a local ring, it suffices to show
  $\overline{f'(\beta)} \ne 0$ in $k_S$
  ($\mathtt{IsLocalRing.}\allowbreak\mathtt{residue\_ne\_zero\_iff\_isUnit}$).

  Since $R \to S$ is étale, the residue field extension $k_R \to k_S$ is
  separable.  By Lemma~\ref{lem:minpoly_map_residue},
  $\bar f = \operatorname{minpoly}_{k_R} \bar\beta$,
  so $\bar f$ is separable
  ($\mathtt{Algebra.IsSeparable.isSeparable}$).
  This means $\bar{f'}(\bar\beta) \ne 0$
  ($\mathtt{Separable.aeval\_derivative\_ne\_zero}$).

  By $\mathtt{hom\_eval_2}$, $\overline{f'(\beta)} = \bar{f'}(\bar\beta) \ne 0$.
\end{proof}

\section{Existence of Generator (Nakayama Argument)}

\begin{theorem}[Finite étale local extensions are monogenic]
  \label{thm:exists_adjoin_eq_top}
  \lean{Monogenic.exists_adjoin_eq_top}
  \leanok
  Let $R \to S$ be a finite étale extension of local rings with faithful
  scalar action.  Then there exists $\beta \in S$ such that $R[\beta] = S$.
\end{theorem}

\begin{proof}
  \leanok
  By the primitive element theorem
  ($\mathtt{Field.exists\_primitive\_element}$),
  there exists $\beta_0 \in k_S$ with $k_R[\beta_0] = k_S$.
  Lift $\beta_0$ to $\beta \in S$ via
  $\mathtt{Ideal.Quotient.mk\_surjective}$.

  It suffices to show $\operatorname{Algebra.adjoin} R \{\beta\} = \top$.
  By Nakayama's lemma ($\mathtt{Submodule.}\allowbreak\mathtt{le\_of\_le\_smul\_of\_le\_jacobson\_bot}$),
  it suffices to show $S \subseteq R[\beta] + \mathfrak{m}_R \cdot S$.

  For any $s \in S$, the residue $\bar s \in k_S = k_R[\beta_0]$, so there exists
  $t_0 \in R[\beta]$ with $\overline{t_0} = \bar s$.  Then
  $s - t_0 \in \mathfrak{m}_S = \mathfrak{m}_R \cdot S$
  (using the étale condition
  $\mathtt{Algebra.FormallyUnramified.map\_maximalIdeal}$).
\end{proof}


\chapter{Converse: Monogenic with Unit Derivative Implies \'Etale}
\label{ch:converse}

\begin{theorem}[Standard étale from monogenic with unit derivative]
  \label{thm:algebra_etale}
  \lean{IsAdjoinRootMonic.algebra_etale}
  \leanok
  Let $f \in R[X]$ be monic with $S \cong R[X]/(f)$
  (via an $\mathtt{IsAdjoinRootMonic}$ structure).
  If $f'(\beta) \in S^\times$ where $\beta$ is the image of $X$,
  then $R \to S$ is étale.
\end{theorem}

\begin{proof}
  \leanok
  Construct a $\mathtt{StandardEtalePair}$ using $(f, f')$:
  the Bézout relation $f' \cdot 1 + f \cdot 0 = (f')^1$ is trivially satisfied.

  The natural map $P.\mathtt{lift} \colon P.\mathtt{Ring} \to S$ sending $X \mapsto \beta$
  is well-defined by $f(\beta) = 0$ and the unit condition on $f'(\beta)$.

  Construct an explicit inverse via $\mathtt{AdjoinRoot.liftAlgHom}$
  composed with the algebra equivalence $S \cong \mathtt{AdjoinRoot}\; f$.
  Verify it is both a left and right inverse by checking on generators,
  concluding bijectivity.

  Hence $S$ is standard étale ($\mathtt{Algebra.IsStandardEtale}$),
  and étale follows.
\end{proof}


\chapter{The Partially \'Etale Case (Lemma 3.1)}
\label{ch:nonetale}

This chapter proves the more general Lemma 3.1: monogenicity from an
étale quotient through a height-one prime. The hypotheses are stronger
on the rings (domains, integrally closed, UFD) but weaker on the map
(only a quotient needs to be étale).

\section{Sub-lemmas}

\begin{lemma}[Height-one primes are principal in UFDs]
  \label{lem:height_one_principal}
  \lean{Monogenic.Ideal.exists_span_singleton_eq_of_prime_of_height_one}
  \leanok
  In a UFD $S$, every prime ideal $\mathfrak{q}$ of height~1 is principal.
\end{lemma}

\begin{proof}
  \leanok
  Since $\mathfrak{q} \ne 0$, it contains an irreducible element $p$.
  Then $\operatorname{span}\{p\}$ is prime and contained in $\mathfrak{q}$.
  If $\operatorname{span}\{p\} \subsetneq \mathfrak{q}$, then
  $\operatorname{height}(\operatorname{span}\{p\}) < \operatorname{height}(\mathfrak{q}) = 1$,
  forcing $\operatorname{span}\{p\} = 0$, contradicting $p \ne 0$.
  Hence $\mathfrak{q} = \operatorname{span}\{p\}$.
\end{proof}

\begin{lemma}[Taylor expansion for polynomial evaluation]
  \label{lem:taylor_expansion}
  \lean{Monogenic.exists_aeval_add_eq}
  \leanok
  For any polynomial $f \in R[X]$ and elements $x, h \in S$, there exists $c \in S$
  such that
  \[
    f(x + h) = f(x) + f'(x) \cdot h + h^2 \cdot c.
  \]
\end{lemma}

\begin{proof}
  \leanok
  Reduce modulo $\operatorname{span}\{h^2\}$ using
  $\mathtt{Polynomial.aeval\_add\_of\_sq\_eq\_zero}$, then lift.
\end{proof}

\begin{lemma}[Maximal ideal decomposition from étale quotient]
  \label{lem:maximalIdeal_eq_sup}
  \lean{Monogenic.maximalIdeal_eq_sup_of_etale_quotient}
  \leanok
  Let $R$ and $S$ be local rings with $S$ a finite $R$-algebra, and let
  $\mathfrak{q} \subset S$ be a prime ideal.  If the quotient map
  $R/(\mathfrak{q} \cap R) \to S/\mathfrak{q}$ is étale, then
  \[
    \mathfrak{m}_S = \mathfrak{q} + \mathfrak{m}_R \cdot S.
  \]
\end{lemma}

\begin{proof}
  \leanok
  The quotient rings are local domains, and the étale condition on
  $R/\mathfrak{p} \to S/\mathfrak{q}$ gives
  $\mathfrak{m}_{R/\mathfrak{p}} \cdot (S/\mathfrak{q}) = \mathfrak{m}_{S/\mathfrak{q}}$
  by $\mathtt{Algebra.FormallyUnramified.map\_maximalIdeal}$.
  Pulling back through the quotient maps, using
  $\operatorname{map}(\operatorname{mk}\; \mathfrak{q})(\mathfrak{m}_S) = \mathfrak{m}_{S/\mathfrak{q}}$
  and the analogous statement for $R$, the equality lifts to
  $\mathfrak{m}_S = \mathfrak{q} + \mathfrak{m}_R S$.
\end{proof}

\begin{lemma}[Quotient adjustment]
  \label{lem:quotient_adjust}
  \lean{Monogenic.Ideal.quotient_adjust}
  \leanok
  Let $\mathfrak{q} \subset S$ be an ideal and $\mathfrak{q} = \operatorname{span}\{q_0\}$.
  Then for any $\beta \in S$, the image of $\beta + q_0$ in $S/\mathfrak{q}$ is
  equal to the image of $\beta$.
\end{lemma}

\begin{proof}
  \leanok
  Since $(\beta + q_0) - \beta = q_0 \in \mathfrak{q}$, the elements are congruent
  modulo $\mathfrak{q}$.
\end{proof}

\section{Nakayama Helpers}


\begin{lemma}[Quotient lifting]
  \label{lem:exists_adjoin_sub_mem}
  \lean{Monogenic.exists_adjoin_sub_mem}
  \leanok
  Let $\mathfrak{q} \subset S$ be an ideal and $\beta \in S$.
  If $\beta$ generates $S/\mathfrak{q}$ over $R/(\mathfrak{q} \cap R)$,
  then for every $s \in S$ there exists $t \in R[\beta]$ with $s - t \in \mathfrak{q}$.
\end{lemma}

\begin{proof}
  \leanok
  Since $(R/\mathfrak{p})[\bar\beta] = S/\mathfrak{q}$, the map
  $\operatorname{aeval}_{\bar\beta}$ is surjective.
  Lift a polynomial preimage of $\bar s$ through
  $\mathtt{Polynomial.map\_surjective}$ to obtain $t \in R[\beta]$
  with $s - t \in \mathfrak{q}$.
\end{proof}

\begin{lemma}[Nakayama argument for quotient generation]
  \label{lem:adjoin_eq_top_of_quotient}
  \lean{Monogenic.adjoin_eq_top_of_quotient}
  \leanok
  \uses{lem:exists_adjoin_sub_mem}
  Let $R, S$ be local rings with $S$ a finite $R$-algebra, and $\mathfrak{q} \subset S$
  a prime ideal.  Let $\beta \in S$ and $\pi \in R[\beta]$.
  Assume:
  \begin{enumerate}
    \item $\bar\beta$ generates $S/\mathfrak{q}$ over $R/(\mathfrak{q} \cap R)$,
    \item $\mathfrak{m}_S = \operatorname{span}\{\pi\} + \mathfrak{m}_R S$.
  \end{enumerate}
  Then $R[\beta] = S$.
\end{lemma}

\begin{proof}
  \leanok
  \uses{lem:exists_adjoin_sub_mem}
  Set $A = R[\beta]$.  The quotient $S/\mathfrak{m}_R S$ is Artinian
  (finite over the field $k_R$), so $\mathfrak{m}_S^n \subseteq \mathfrak{m}_R S$
  for some $n$.

  \textbf{Power bound.} By induction,
  $\mathfrak{m}_S^k \subseteq \operatorname{span}\{\pi^k\} + \mathfrak{m}_R S$,
  using the hypothesis on $\mathfrak{m}_S$.

  \textbf{Iterative reduction.} For $x \in \mathfrak{q}$, by repeated application
  of quotient lifting (Lemma~\ref{lem:exists_adjoin_sub_mem}) and the power bound,
  we find $a \in A$ with $x - a \in \operatorname{span}\{\pi^n\} + \mathfrak{m}_R S
  \subseteq \mathfrak{m}_R S$.

  Hence $\mathfrak{q} \subseteq A + \mathfrak{m}_R S$ as $R$-submodules.
  Combined with the quotient lifting for general elements,
  $S \subseteq A + \mathfrak{m}_R S$.
  By Nakayama's lemma, $A = S$.
\end{proof}

\section{Principal Adjustment Lemma}

\begin{lemma}[Adjustment by principal generator]
  \label{lem:exists_isAdjoinRootMonic_of_principal_adjust}
  \lean{Monogenic.exists_isAdjoinRootMonic_of_principal_adjust}
  \leanok
  \uses{lem:taylor_expansion, lem:adjoin_eq_top_of_quotient, thm:mkOfAdjoinEqTop, lem:quotient_adjust}
  Let $(R, \mathfrak{m}_R)$ and $(S, \mathfrak{m}_S)$ be local domains with
  $R$ integrally closed and $S$ a finite $R$-algebra with faithful scalar action.
  Let $\mathfrak{q} \subset S$ be a prime ideal with $\mathfrak{q} = (q_0)$,
  and let $\beta \in S$, $f_1 \in R[X]$, $a \in S$ satisfy:
  \begin{enumerate}
    \item $f_1(\beta) = q_0 \cdot a$ with $a \in \mathfrak{m}_S$,
    \item $f_1'(\beta) \notin \mathfrak{m}_S$,
    \item $\mathfrak{m}_S = \mathfrak{q} + \mathfrak{m}_R S$,
    \item $\bar\beta$ generates $S/\mathfrak{q}$ over $R/(\mathfrak{q} \cap R)$.
  \end{enumerate}
  Then there exist a monic $f \in R[X]$ and an $\mathtt{IsAdjoinRootMonic}$ structure
  for $S$ over $f$.
\end{lemma}

\begin{proof}
  \leanok
  \uses{lem:taylor_expansion, lem:adjoin_eq_top_of_quotient, thm:mkOfAdjoinEqTop, lem:quotient_adjust}
  Set $\beta' = \beta + q_0$.
  By Lemma~\ref{lem:taylor_expansion},
  $f_1(\beta') = q_0 \cdot (a + f_1'(\beta) + q_0 b)$ for some $b \in S$.
  Since $f_1'(\beta) \notin \mathfrak{m}_S$
  and $a, q_0 b \in \mathfrak{m}_S$ (the former by hypothesis), the cofactor is a unit.
  Hence $\operatorname{span}\{f_1(\beta')\} = \mathfrak{q}$, and
  $\mathfrak{m}_S = \operatorname{span}\{f_1(\beta')\} + \mathfrak{m}_R S$.

  Since $\beta' \equiv \beta \pmod{\mathfrak{q}}$ (Lemma~\ref{lem:quotient_adjust}),
  it still generates the quotient, so Lemma~\ref{lem:adjoin_eq_top_of_quotient}
  gives $R[\beta'] = S$.
  Finally, Theorem~\ref{thm:mkOfAdjoinEqTop} provides the
  $\mathtt{IsAdjoinRootMonic}$ structure for $\operatorname{minpoly}_R \beta'$.
\end{proof}

\section{Main Theorem: Lemma 3.1}

\begin{theorem}[Monogenicity from étale height-one quotient]
  \label{thm:exists_isAdjoinRootMonic_of_quotientMap_etale}
  \lean{Monogenic.exists_isAdjoinRootMonic_of_quotientMap_etale}
  \leanok
  \uses{thm:exists_adjoin_eq_top, thm:mkOfAdjoinEqTop, lem:height_one_principal,
    lem:maximalIdeal_eq_sup, lem:adjoin_eq_top_of_quotient,
    thm:isUnit_aeval_derivative_minpoly,
    lem:exists_isAdjoinRootMonic_of_principal_adjust}
  Let $(R, \mathfrak{m}_R)$ and $(S, \mathfrak{m}_S)$ be local domains with
  $R$ integrally closed and $S$ a UFD.
  Let $R \to S$ be a finite injective ring homomorphism.
  If there exists a prime ideal $\mathfrak{q} \subset S$ of height~1 such that
  the induced map $R/(\mathfrak{q} \cap R) \to S/\mathfrak{q}$ is étale, then
  there exist a monic $f \in R[X]$ and an $\mathtt{IsAdjoinRootMonic}$ structure
  for $S$ over $f$.
\end{theorem}

\begin{proof}
  \leanok
  \uses{thm:exists_adjoin_eq_top, thm:mkOfAdjoinEqTop, lem:height_one_principal,
    lem:maximalIdeal_eq_sup, lem:adjoin_eq_top_of_quotient,
    thm:isUnit_aeval_derivative_minpoly,
    lem:exists_isAdjoinRootMonic_of_principal_adjust}
  \textbf{Case 1: $R \to S$ is already étale.}
  Apply Theorem~\ref{thm:exists_adjoin_eq_top} to find $\beta$ with $R[\beta] = S$,
  then Theorem~\ref{thm:mkOfAdjoinEqTop} for the isomorphism.

  \textbf{Case 2: $R \to S$ is not étale.}
  Apply Theorem~\ref{thm:exists_adjoin_eq_top} to the étale quotient
  $R_0 = R/\mathfrak{p} \to S_0 = S/\mathfrak{q}$ to get $\beta_0 \in S_0$ with
  $R_0[\beta_0] = S_0$.  Lift to $\beta \in S$.

  By Lemma~\ref{lem:height_one_principal}, $\mathfrak{q} = (q_0)$ for some $q_0 \in S$.
  Lift the minimal polynomial $f_0 = \operatorname{minpoly}_{R_0} \beta_0$
  to a monic $f_1 \in R[X]$.
  By Lemma~\ref{lem:maximalIdeal_eq_sup},
  $\mathfrak{m}_S = \mathfrak{q} + \mathfrak{m}_R S$.

  \textbf{Sub-case 2a:} If $f_1(\beta)$ generates $\mathfrak{m}_S$ modulo
  $\mathfrak{m}_R S$, apply Lemma~\ref{lem:adjoin_eq_top_of_quotient} directly
  to conclude $R[\beta] = S$.

  \textbf{Sub-case 2b:} Otherwise, since $f_1(\beta) \in \mathfrak{q} = (q_0)$,
  write $f_1(\beta) = q_0 \cdot a$.  If $a$ were a unit, then
  $\operatorname{span}\{f_1(\beta)\} = \mathfrak{q}$ and Sub-case~2a would apply;
  so $a \in \mathfrak{m}_S$.
  The derivative condition $f_1'(\beta) \notin \mathfrak{m}_S$ follows from
  Theorem~\ref{thm:isUnit_aeval_derivative_minpoly} applied to the étale quotient.
  Now apply Lemma~\ref{lem:exists_isAdjoinRootMonic_of_principal_adjust}
  with $q_0$, $\beta$, $f_1$, and $a$.
\end{proof}

% Removed: exists_adjoin_eq_top_of_quotientMap_etale (trivial corollary, not formalized separately)


\chapter{Plan for Mathlib PRs}
\label{ch:mathlib_plan}

This chapter outlines a strategy for upstreaming the formalized results into
Mathlib via a sequence of small, self-contained pull requests.
Each PR should be independently reviewable and useful.

Throughout, we reference Mathlib file paths as of the current toolchain.
The key Mathlib files involved are:
\begin{itemize}
  \item \texttt{Mathlib/RingTheory/IsAdjoinRoot.lean} ---
    defines \texttt{IsAdjoinRoot}, \texttt{IsAdjoinRootMonic},
    and \texttt{finrank\_quotient\_span\_eq\_natDegree'}.
    Imports \texttt{FieldTheory.Minpoly.IsIntegrallyClosed}
    and \texttt{RingTheory.PowerBasis}.
    Does \emph{not} transitively import \texttt{LinearAlgebra.Charpoly.Basic}.
  \item \texttt{Mathlib/RingTheory/Etale/Basic.lean} ---
    defines \texttt{Algebra.Etale}.
  \item \texttt{Mathlib/RingTheory/Etale/StandardEtale.lean} ---
    defines \texttt{StandardEtalePair}, \texttt{StandardEtalePresentation},
    and \texttt{Algebra.IsStandardEtale}.
  \item \texttt{Mathlib/RingTheory/Unramified/LocalRing.lean} ---
    contains \texttt{FormallyUnramified.map\_maximalIdeal}.
  \item \texttt{Mathlib/RingTheory/LocalRing/ResidueField/} ---
    \texttt{Defs.lean} (defines \texttt{residue}, \texttt{ResidueField}),
    \texttt{Basic.lean} (surjectivity, kernel),
    \texttt{Instances.lean} (separability instances).
  \item \texttt{Mathlib/RingTheory/LocalRing/Quotient.lean} ---
    contains \texttt{finrank\_quotient\_map} and
    \texttt{exists\_maximalIdeal\_pow\_le\_of\_isArtinianRing\_quotient}.
  \item \texttt{Mathlib/RingTheory/LocalRing/Module.lean} ---
    contains \texttt{Module.free\_of\_flat\_of\_isLocalRing}.
  \item \texttt{Mathlib/RingTheory/Ideal/Height.lean} ---
    defines \texttt{Ideal.height}.
  \item \texttt{Mathlib/FieldTheory/PrimitiveElement.lean} ---
    contains \texttt{Field.exists\_primitive\_element}.
  \item \texttt{Mathlib/RingTheory/Nakayama.lean} ---
    contains \texttt{Submodule.le\_of\_le\_smul\_of\_le\_}\allowbreak\texttt{jacobson\_bot}.
  \item \texttt{Mathlib/Algebra/Polynomial/Taylor.lean} ---
    contains \texttt{aeval\_add\_of\_sq\_eq\_zero}.
\end{itemize}


%% ===== PR 1 =====
\section{PR 1: Quotient Isomorphism without Integrally Closed}
\label{sec:pr1}

\subsection*{Declarations}

\begin{itemize}
  \item \texttt{minpoly.natDegree\_le'}: For a finite free extension with
    $R$ nontrivial, $\deg(\operatorname{minpoly}_R \alpha) \le \operatorname{finrank}_R S$.
    Generalizes the existing \texttt{minpoly.natDegree\_le} in
    \texttt{FieldTheory/IntermediateField/Adjoin/Basic.lean}
    (which requires $K, L$ to be fields with \texttt{FiniteDimensional}).
  \item \texttt{IsAdjoinRootMonic.mkOfAdjoinEqTop'}: If $R[\alpha] = S$ and $S$ is
    finite free over $R$, then $S$ admits an \texttt{IsAdjoinRootMonic}
    structure for $\operatorname{minpoly}_R \alpha$.
\end{itemize}

\subsection*{File placement}

\textbf{Option A (preferred):} Both declarations go into
\texttt{Mathlib/RingTheory/IsAdjoinRoot.lean}.
This is where \texttt{IsAdjoinRootMonic}, \texttt{finrank\_quotient\_span\_eq\_natDegree'},
and the existing adjoin-root machinery live.

\textbf{Import issue:} \texttt{natDegree\_le'} uses
\texttt{LinearMap.aeval\_self\_charpoly} and \texttt{charpoly\_natDegree}
from \texttt{Mathlib.LinearAlgebra.Charpoly.Basic}, which is \emph{not}
a transitive import of \texttt{IsAdjoinRoot.lean}.
Adding this import may be acceptable since the file already has heavy imports
(\texttt{PowerBasis}, \texttt{Minpoly.IsIntegrallyClosed}), but reviewers
may prefer a separate file.

\textbf{Option B:} Put \texttt{natDegree\_le'} in a new file
\texttt{Mathlib/FieldTheory/Minpoly/Basic.lean} (or
\texttt{Mathlib/RingTheory/Polynomial/Minpoly/Degree.lean}),
then import it from \texttt{IsAdjoinRoot.lean} for \texttt{mkOfAdjoinEqTop'}.

\subsection*{Namespace and naming}

\begin{itemize}
  \item \texttt{minpoly.natDegree\_le'}: Currently uses \texttt{\_root\_} prefix
    to escape the \texttt{Monogenic} namespace. In Mathlib, consider renaming to
    \texttt{minpoly.natDegree\_le\_finrank} or similar,
    and whether to replace the existing field-level
    \texttt{minpoly.natDegree\_le} or keep both.
  \item \texttt{IsAdjoinRootMonic.mkOfAdjoinEqTop'}: Already in the correct
    namespace (matching the Mathlib original \texttt{IsAdjoinRootMonic.mkOfAdjoinEqTop}).
    For Mathlib, the prime suffix should be dropped; the existing Mathlib version
    requires \texttt{IsIntegrallyClosed}, so this would replace it.
\end{itemize}

\subsection*{Required imports (for the declarations)}
\begin{itemize}
  \item \texttt{Mathlib.LinearAlgebra.Charpoly.Basic}
    (for \texttt{aeval\_self\_charpoly}, \texttt{charpoly\_natDegree},
     \texttt{charpoly\_monic}).
  \item \texttt{Mathlib.RingTheory.OrzechProperty}
    (for \texttt{injective\_of\_surjective\_endomorphism};
     likely already transitive via \texttt{IsAdjoinRoot}).
  \item \texttt{Mathlib.RingTheory.IsAdjoinRoot}
    (for \texttt{IsAdjoinRootMonic}, \texttt{AdjoinRoot.liftAlgHom},
     \texttt{finrank\_quotient\_span\_eq\_natDegree'}).
\end{itemize}

\subsection*{Proof adjustments for Mathlib style}
\begin{itemize}
  \item The current proof of \texttt{natDegree\_le'} uses \texttt{classical}
    and local \texttt{let} bindings for basis and matrix. This should be
    rewritten to avoid \texttt{classical} if possible (use \texttt{Decidable}
    instances or restructure).
  \item Drop the \texttt{Monogenic} namespace wrapper.
  \item Add docstrings in Mathlib format.
  \item The current proof uses
    \texttt{Module.finrank\_eq\_card\_chooseBasisIndex};
    check this is still the canonical spelling.
\end{itemize}

\subsection*{Potential issues}
\begin{itemize}
  \item Import weight: adding \texttt{Charpoly.Basic} to the import tree of
    \texttt{IsAdjoinRoot.lean} may increase compile times. Measure before submitting.
  \item The existing \texttt{minpoly.natDegree\_le} for fields is proved via
    \texttt{IntermediateField.adjoin.finrank}. Reviewers may ask if the
    field version should be deduced from the ring version or kept separate.
\end{itemize}


%% ===== PR 2 =====
\section{PR 2: Residue Field Properties for \'Etale Local Extensions}
\label{sec:pr2}

\subsection*{Declarations}

\begin{itemize}
  \item \texttt{adjoin\_residue\_eq\_top\_of\_adjoin\_eq\_top}:
    If $R[\beta] = S$ then $k_R[\bar\beta] = k_S$.
  \item \texttt{finrank\_eq\_finrank\_residueField}:
    $\operatorname{finrank}_R S = \operatorname{finrank}_{k_R} k_S$
    for finite étale local extensions.
  \item \texttt{minpoly\_map\_residue}: For a finite étale local extension
    with $R[\beta] = S$:
    $(\operatorname{minpoly}_R \beta) \bmod \mathfrak{m}_R
    = \operatorname{minpoly}_{k_R} \bar\beta$.
  \item \texttt{isUnit\_aeval\_derivative\_minpoly\_of\_adjoin\_eq\_top}:
    For a finite étale local extension with $R[\beta] = S$,
    $f'(\beta) \in S^\times$ where $f = \operatorname{minpoly}_R \beta$.
\end{itemize}

\subsection*{File placement}

\textbf{Preferred:} A new file \texttt{Mathlib/RingTheory/LocalRing/ResidueField/Minpoly.lean} or \texttt{Mathlib/RingTheory/Etale/Monogenic.lean} that collects these cohesive residue field properties for étale local extensions.

\subsection*{Namespace}

\texttt{IsLocalRing} or \texttt{Algebra.Etale} namespace.

\subsection*{Required imports}
\begin{itemize}
  \item \texttt{Mathlib.RingTheory.LocalRing.ResidueField.Defs}
    (for \texttt{residue}, \texttt{ResidueField}).
  \item \texttt{Mathlib.RingTheory.Unramified.LocalRing}
    (for \texttt{FormallyUnramified.}\allowbreak\texttt{map\_maximalIdeal},
     used by \texttt{finrank\_eq\_finrank\_residueField}).
  \item \texttt{Mathlib.RingTheory.Smooth.Flat}
    (for \texttt{Algebra.Smooth.flat}, used in the étale $\Rightarrow$ free chain).
  \item \texttt{Mathlib.RingTheory.LocalRing.Module}
    (for \texttt{free\_of\_flat\_of\_isLocalRing}).
  \item \texttt{Mathlib.RingTheory.LocalRing.Quotient}
    (for \texttt{finrank\_quotient\_map}).
  \item \texttt{IsAdjoinRootMonic.mkOfAdjoinEqTop'} from PR~1
    (used in the degree chain: $\deg f = \operatorname{finrank}_R S$
     via \texttt{.finrank}).
  \item \texttt{Mathlib.FieldTheory.IntermediateField.Adjoin.Basic}
    (for \texttt{IntermediateField.adjoin.finrank},
     \texttt{adjoin\_simple\_toSubalgebra\_}\allowbreak\texttt{of\_isAlgebraic}).
  \item \texttt{Mathlib.FieldTheory.Separable}
    (for \texttt{Algebra.IsSeparable.isSeparable},
     \texttt{Separable.aeval\_derivative\_ne\_zero}).
\end{itemize}

\subsection*{Proof structure}
\begin{itemize}
  \item \textbf{Residue generator \& rank:} \texttt{finrank\_eq\_finrank\_residueField} uses étale $\Rightarrow$ smooth $\Rightarrow$ flat $\Rightarrow$ free.
  \item \textbf{Minimal polynomial descent:} \texttt{minpoly\_map\_residue} compares degrees:
    $g \mid \bar f$, both monic, same degree $\Rightarrow$ equal.
    The degree chain uses \texttt{mkOfAdjoinEqTop'.finrank} (PR~1),
    \texttt{finrank\_eq\_finrank\_residueField}, and
    \texttt{IntermediateField.adjoin.finrank} (Mathlib).
  \item \textbf{Unit derivative:} Reduce to residue field via \texttt{minpoly\_map\_residue},
    apply separability of $k_R \to k_S$ (from étale), conclude
    $\bar{f'}(\bar\beta) \ne 0$, lift back to unit via
    \texttt{IsLocalRing.residue\_ne\_zero\_iff\_isUnit}.
\end{itemize}

\subsection*{Potential issues}
\begin{itemize}
  \item \texttt{finrank\_eq\_finrank\_residueField} uses the étale
    $\Rightarrow$ smooth $\Rightarrow$ flat chain. Consider whether to state it more generally for
    ``finite free local extensions where $\mathfrak{m}_R S = \mathfrak{m}_S$''.
  \item The étale hypothesis is used only to get
    $\operatorname{finrank}_R S = \operatorname{finrank}_{k_R} k_S$.
    If this equality holds more generally, the \texttt{minpoly\_map\_residue} lemma
    could be stated without \texttt{Algebra.Etale}. Reviewers may ask for
    this generalization.
  \item The hypothesis chain étale $\Rightarrow$ separable residue field
    should be automatic, but verify the instance path works:
    \texttt{Algebra.Etale R S} $\Rightarrow$
    \texttt{Algebra.FormallyUnramified R S} $\Rightarrow$
    \texttt{Algebra.IsSeparable kR kS}.
  \item Uses \texttt{eq\_of\_monic\_of\_dvd\_of\_natDegree\_le} --- verify
    this is still the canonical Mathlib name.
\end{itemize}


%% ===== PR 3 =====
\section{PR 3: Finite \'Etale Local Extensions are Monogenic}
\label{sec:pr3}

\subsection*{Declarations}

\begin{itemize}
  \item \texttt{exists\_adjoin\_eq\_top}: For a finite étale local extension
    $R \to S$ with faithful scalar action,
    $\exists \beta \in S,\; R[\beta] = S$.
\end{itemize}

\subsection*{File placement}

\textbf{Preferred:} New file \texttt{Mathlib/RingTheory/Etale/Monogenic.lean}
or \texttt{Mathlib/RingTheory/Etale/LocalRing.lean}.
This is a natural home for étale-specific results about local rings.

\subsection*{Namespace}

\texttt{Algebra.Etale.exists\_adjoin\_eq\_top} or just in the root namespace
with appropriate variable scoping.

\subsection*{Required imports}
\begin{itemize}
  \item \texttt{Mathlib.FieldTheory.PrimitiveElement}
    (for \texttt{Field.exists\_primitive\_element}).
  \item \texttt{Mathlib.RingTheory.Unramified.LocalRing}
    (for \texttt{FormallyUnramified.}\allowbreak\texttt{map\_maximalIdeal}).
  \item \texttt{Mathlib.RingTheory.Nakayama}
    (for \texttt{Submodule.le\_of\_le\_smul\_of\_le\_}\allowbreak\texttt{jacobson\_bot}).
  \item \texttt{Mathlib.RingTheory.LocalRing.ResidueField.Defs}
    (for \texttt{residue}, \texttt{ResidueField}).
\end{itemize}

\subsection*{Proof structure}

Primitive element theorem gives $\beta_0 \in k_S$ with $k_R[\beta_0] = k_S$.
Lift to $\beta \in S$.
Show $S \subseteq R[\beta] + \mathfrak{m}_R S$: for each $s \in S$,
lift a polynomial preimage of $\bar s$ through
$\mathtt{Polynomial.map\_surjective}$ to find $t \in R[\beta]$
with $s - t \in \mathfrak{m}_S = \mathfrak{m}_R S$.
Conclude by Nakayama.

\subsection*{Potential issues}
\begin{itemize}
  \item This PR is \emph{independent} of PRs 1 and 2.
    It does not need \texttt{mkOfAdjoinEqTop'} or
    \texttt{minpoly\_map\_residue}; it only produces a generator $\beta$,
    not an isomorphism $S \cong R[X]/(f)$.
  \item The proof uses \texttt{IntermediateField.adjoin\_simple\_toSubalgebra\_of\_isAlgebraic}
    to convert between \texttt{Algebra.adjoin} and \texttt{IntermediateField.adjoin}
    in the residue field. This is a common pattern but sometimes fragile with
    universe issues.
  \item The \texttt{FaithfulSMul R S} hypothesis is equivalent to injectivity
    of $\varphi$. Reviewers may prefer \texttt{IsLocalHom (algebraMap R S)}
    or just \texttt{Function.Injective (algebraMap R S)}.
    Note \texttt{faithfulSMul\_iff\_algebraMap\_injective} exists in Mathlib.
\end{itemize}


%% ===== PR 4 =====
\section{PR 4: Converse --- Standard \'Etale from Unit Derivative}
\label{sec:pr4}

\subsection*{Declarations}

\begin{itemize}
  \item \texttt{IsAdjoinRootMonic.algebra\_etale}: If $S \cong R[X]/(f)$
    (via \texttt{IsAdjoinRootMonic}) and $f'(\beta) \in S^\times$,
    then $R \to S$ is étale.
\end{itemize}

\subsection*{File placement}

\textbf{Preferred:} \texttt{Mathlib/RingTheory/Etale/StandardEtale.lean},
directly after the definition of \texttt{StandardEtalePair} and
\texttt{Algebra.IsStandardEtale}. This is where the standard étale
machinery lives and the proof constructs a \texttt{StandardEtalePair}.

\subsection*{Namespace}

\texttt{IsAdjoinRootMonic.algebra\_etale} (already correctly namespaced).

\subsection*{Required imports}
\begin{itemize}
  \item \texttt{Mathlib.RingTheory.Etale.StandardEtale}
    (for \texttt{StandardEtalePair}, \texttt{HasMap}, \texttt{lift},
     \texttt{hom\_ext}, \texttt{Algebra.IsStandardEtale}).
  \item \texttt{Mathlib.RingTheory.IsAdjoinRoot}
    (for \texttt{IsAdjoinRootMonic}, \texttt{adjoinRootAlgEquiv},
     \texttt{map\_surjective}).
\end{itemize}

Both of these should already be in the import tree of
\texttt{StandardEtale.lean}, so no new imports are needed if placed there.

\subsection*{Proof structure}

Construct \texttt{StandardEtalePair R} with $f, f', 1, 0, 1$ and the
trivial Bézout relation $f' \cdot 1 + f \cdot 0 = (f')^1$.
Build an explicit inverse of \texttt{P.lift} by composing
\texttt{AdjoinRoot.liftAlgHom} with \texttt{adjoinRootAlgEquiv.symm}.
Verify left/right inverse on generators via \texttt{P.hom\_ext}.

\subsection*{Potential issues}
\begin{itemize}
  \item This theorem essentially says ``standard étale presentations arise
    from adjoin-root-monic with unit derivative.'' Mathlib may already have
    something close or the reviewers may want to restructure. Check for
    overlap with existing material in \texttt{StandardEtale.lean}.
  \item The proof uses \texttt{AdjoinRoot.liftAlgHom} which was recently
    introduced as a replacement for the deprecated \texttt{AdjoinRoot.liftHom}.
    Verify this is still the current API.
\end{itemize}


%% ===== PR 5 =====
\section{PR 5: Nakayama Helpers and Partially \'Etale Case (Lemma 3.1)}
\label{sec:pr5}

This is the most complex PR. It should likely be split into several sub-PRs.

\subsection*{Sub-PR 5a: Height-One Primes in UFDs are Principal}

\textbf{Declaration:}
\texttt{Ideal.exists\_span\_singleton\_eq\_of\_prime\_of\_height\_one}

\textbf{File:} \texttt{Mathlib/RingTheory/Ideal/Height.lean} (append to existing file).

\textbf{Imports:} Only needs \texttt{Ideal.height}, \texttt{IsPrime},
\texttt{UniqueFactorizationMonoid}. All already in that file's import tree.

\textbf{Notes:} Self-contained, no dependencies on other PRs.
This is a basic fact about UFDs that should be independently useful.
The proof uses \texttt{IsPrime.exists\_mem\_prime\_of\_ne\_bot},
\texttt{height\_strict\_mono\_of\_is\_prime}, and
\texttt{primeHeight\_eq\_zero\_iff}.

\subsection*{Sub-PR 5b: Taylor Expansion for Polynomial Evaluation}

\textbf{Declaration:}
\texttt{exists\_aeval\_add\_eq}: $f(x+h) = f(x) + f'(x)h + h^2 c$.

\textbf{File:} \texttt{Mathlib/Algebra/Polynomial/Taylor.lean}
(append to existing file, which already has
\texttt{aeval\_add\_of\_sq\_eq\_zero}).

\textbf{Imports:} Only needs \texttt{Polynomial.Taylor} (already there)
and \texttt{Ideal.Quotient} basics.

\textbf{Notes:} Self-contained. The proof reduces mod $\langle h^2 \rangle$
using \texttt{aeval\_add\_of\_sq\_eq\_zero}, then lifts via
\texttt{Ideal.mem\_span\_singleton}. Very clean.

\subsection*{Sub-PR 5c: Maximal Ideal Decomposition and Quotient Lifting}

\textbf{Declarations:}
\begin{itemize}
  \item \texttt{maximalIdeal\_eq\_sup\_of\_etale\_quotient}:
    If $R/\mathfrak{p} \to S/\mathfrak{q}$ is étale, then
    $\mathfrak{m}_S = \mathfrak{q} + \mathfrak{m}_R S$.
  \item \texttt{exists\_adjoin\_sub\_mem}: Quotient lifting.
  \item \texttt{adjoin\_eq\_top\_of\_quotient}: Nakayama with quotient generation.
\end{itemize}

\textbf{File:} New file, perhaps
\texttt{Mathlib/RingTheory/LocalRing/Quotient/Etale.lean} or
\texttt{Mathlib/RingTheory/Etale/LocalRing.lean}.

\textbf{Imports:}
\begin{itemize}
  \item \texttt{Mathlib.RingTheory.Unramified.LocalRing}
    (for \texttt{FormallyUnramified.}\allowbreak\texttt{map\_maximalIdeal}).
  \item \texttt{Mathlib.RingTheory.RingHom.Etale}
    (for \texttt{RingHom.Etale}).
  \item \texttt{Mathlib.RingTheory.LocalRing.Quotient}
    (for \texttt{exists\_maximalIdeal\_pow\_le\_of\_}\allowbreak\texttt{isArtinianRing\_quotient}).
  \item \texttt{Mathlib.RingTheory.Nakayama}.
\end{itemize}

\textbf{Notes:}
\begin{itemize}
  \item \texttt{maximalIdeal\_eq\_sup\_of\_}\allowbreak\texttt{etale\_quotient} has a somewhat
    long proof involving quotient maps and the commutative square
    $\operatorname{map}(\operatorname{mk}\;\mathfrak{q})$.
    It uses \texttt{Ideal.map\_eq\_iff\_sup\_ker\_}\allowbreak\texttt{eq\_of\_surjective}
    which may be fragile.
  \item \texttt{adjoin\_eq\_top\_of\_quotient} is the most intricate proof
    in the project (Artinian descent + iterative reduction + Nakayama).
    It is already decomposed into helper lemmas:
    \texttt{Ideal.pow\_le\_span\_pow\_sup} (pure ideal arithmetic)
    and \texttt{exists\_sub\_mem\_adjoin\_of\_pow} (iterative approximation).
  \item \texttt{exists\_adjoin\_sub\_mem} lifts polynomials through
    the quotient map via \texttt{Polynomial.map\_surjective}.
\end{itemize}

\subsection*{Sub-PR 5d: Main Theorem (Lemma 3.1)}

\textbf{Declarations:}
\begin{itemize}
  \item \texttt{exists\_isAdjoinRootMonic\_of\_principal\_adjust}:
    The principal adjustment sub-case (Case 2b), extracted as a
    standalone lemma.
  \item \texttt{exists\_isAdjoinRootMonic\_of\_quotientMap\_etale}:
    The main theorem.
\end{itemize}

\textbf{File:} Same as 5c, or a dedicated file
\texttt{Mathlib/RingTheory/Etale/Monogenic.lean}.

\textbf{Dependencies:} PRs 1, 2, 3, and sub-PRs 5a--5c.

\textbf{Imports (beyond 5c):}
\begin{itemize}
  \item \texttt{Mathlib.RingTheory.Ideal.Height} (for height-one primes, 5a).
  \item \texttt{Mathlib.Algebra.Polynomial.Taylor} (for Taylor expansion, 5b).
  \item \texttt{Mathlib.Algebra.Polynomial.Lifts}
    (for \texttt{lifts\_and\_degree\_eq\_and\_monic}, used to lift
     $f_0 \in R_0[X]$ to monic $f_1 \in R[X]$).
  \item PR~1 (\texttt{mkOfAdjoinEqTop'}, used to build the isomorphism).
  \item PR~2 (\texttt{isUnit\_aeval\_derivative\_minpoly\_of\_adjoin\_eq\_top}, for the
     derivative non-vanishing in Case 2b).
  \item PR~3 (\texttt{exists\_adjoin\_eq\_top}, applied to the étale quotient).
\end{itemize}

\textbf{Proof structure:} Case split on whether $R \to S$ is already étale.
Case 1 uses PR~3 + PR~1 directly.
Case 2 lifts a generator from the étale quotient, then either
(2a) applies the Nakayama argument directly, or
(2b) extracts $a$ with $f_1(\beta) = q_0 \cdot a$ and shows $a \in \mathfrak{m}_S$
(if $a$ were a unit, Sub-case~2a would apply), then delegates to
\texttt{exists\_isAdjoinRootMonic\_of\_principal\_adjust}.

\textbf{Notes:}
\begin{itemize}
  \item The hypotheses include \texttt{IsIntegrallyClosed R} and
    \texttt{UniqueFactorizationMonoid S}. These are strong; reviewers may
    ask whether they can be weakened.
  \item An alternative formulation with an explicit ring homomorphism
    (transferring between $\operatorname{Algebra.adjoin}\;R$ and
    $\operatorname{Algebra.adjoin}\;\varphi(R)$) could be added as a
    trivial corollary if needed.
\end{itemize}


%% ===== Dependency Graph =====
\section{Dependency Graph Summary}

The PR dependency structure is:
\[
  \begin{array}{ccccc}
    \text{PR 1} & & \text{PR 3} & & \text{PR 4} \\
    \downarrow & & \downarrow & & \\
    \text{PR 2} & \to & \text{PR 5} & & \\
  \end{array}
\]

PRs 1, 3, and 4 can be submitted independently and in parallel.
PR 2 depends on PR 1.
PR 5 depends on PRs 1, 2, and 3.
Sub-PRs 5a and 5b are fully independent and can be submitted first.
